\section{案例拆解：从任务描述到可执行轨迹}

\subsection{以 ``预订符合约束的餐厅'' 为例}
讲座中多次以餐厅检索类任务说明 Agent 的执行难点。这类任务看似简单，但在真实网页里通常包含多重约束：地点、价格、评分、营业时间、菜系偏好以及下单流程。系统如果没有显式状态管理，很容易在中途忘记条件。

\begin{importantbox}{任务拆解的推荐结构}
\begin{enumerate}
    \item \textbf{约束抽取}：把自然语言需求转为结构化字段；
    \item \textbf{候选检索}：在站内或跨站生成候选集合；
    \item \textbf{证据核验}：逐项核验硬约束（如价格上限）；
    \item \textbf{动作执行}：进入下单流程并完成关键字段填写；
    \item \textbf{终局确认}：通过页面状态与 verifier 双重确认任务完成。
\end{enumerate}
\end{importantbox}

\begin{knowledgebox}{为什么 ``结构化约束'' 很关键}
如果不把约束结构化，Agent 往往会在长链路中丢失关键信息，只保留模糊偏好。把约束显式存成键值对后，可在每一步执行前做一致性检查，显著降低 ``后期才发现不满足条件'' 的返工成本。
\end{knowledgebox}

\subsection{轨迹级质量控制}
在该任务中，单步正确并不代表轨迹正确。例如模型可能先选到看起来不错的餐厅，但价格或可用时间不满足要求。轨迹级质量控制要检查 ``目标一致性''，而不是只检查 ``动作可执行性''。

\begin{table}[H]
\centering
\small
\begin{tabular}{p{2.8cm}p{3.5cm}p{4.4cm}}
\toprule
\textbf{检查层} & \textbf{检查对象} & \textbf{常见错误} \\
\midrule
字段一致性 & 价格、时间、位置、评分 & 字段遗漏、单位混淆、约束冲突 \\
页面一致性 & 当前页面是否包含目标证据 & 元素识别错误、读错卡片信息 \\
流程一致性 & 是否进入正确下单路径 & 错站点、错入口、提前停止 \\
终局一致性 & 任务结果是否与原需求一致 & 局部成功但全局不满足 \\
\bottomrule
\end{tabular}
\caption{餐厅检索任务的轨迹级质量控制}
\end{table}

\begin{figure}[H]
\centering
\includegraphics[width=0.82\textwidth]{frame-07-search-eval.jpg}
\caption{视频约 00:34:39：在候选轨迹间比较后再继续扩展，体现 ``先评估后执行'' 的策略}
\end{figure}

\subsection{可复用的执行模板}
为了减少每次任务都从零探索，可维护 ``任务模板库''。模板不是固定脚本，而是 ``可约束的流程骨架''：定义阶段顺序、关键检查点与回滚策略。这样既保留灵活性，又能显著降低长任务方差。

\begin{warningbox}{模板化的边界}
模板过刚会损害泛化，模板过松又无法约束错误。合理做法是将模板用于高风险步骤（支付、身份、确认页），而将探索留给低风险步骤（候选检索与排序）。
\end{warningbox}

\subsection{本章小结}
复杂网页任务的关键不在于 ``让模型更会猜''，而在于 ``让轨迹更可控''。结构化约束、轨迹级检查和可复用模板是把演示系统变成稳定系统的三件基础工作。

